\SourceLeaf{289}%
\begin{center}
\begin{tabular}{lr}
Chasseral-Suchet&\(+0''04536\)\\
Chasseral-Berra&\(-0{,}00966\)\\
Suchet-Chasseral&\(+0{,}06221\)\\
Suchet-Berra&\(+0{,}01014\)\\
Berra-Suchet&\(-0{,}04717\)\\
Berra-Chasseral&\(-0{,}06039\)
\end{tabular}
\end{center}
so here too they have no effect on the calculation, which in the work cited is carried out to tenths of a second.
\Article{15}
The rules given solve the problem treated in Arts. 12 and 13 with more accuracy than its application requires. Still, it is worth carrying the accuracy one order further in a special case, and doing so is even necessary to complete, on a consistent basis, an investigation to be presented later. This special case is subject to the condition that the line \(N\) meets the standard parallel between \(F\) and \(G\), at a point \(H\). In this case it is better to move the origin of \(x\) from \(F\), where it was placed above, to \(H\). As a result, when \(l\) and \(m\) are expanded in series in ascending powers of \(x\), the first and second terms of the former series and the second and third terms of the latter drop out. The series therefore take the following form
\[
\begin{aligned}
l&=\lambda xx+\lambda'x^3+\text{etc.}\\
m&=1+\mu x^3+\mu'x^4+\text{etc.}
\end{aligned}
\]
\indent For our purpose only one coefficient in these series, \(\lambda\), will be needed. The following expression for it follows readily from the formula for \(\log m\) given in Art. 9, combined with the expansions in Art. 13
\[
\lambda=-\frac{2ee\cos P\sin P\sin\chi\cos\chi^2}{\cos\varphi\cos\theta}
\]
Here \(e,P,\varphi,\theta\) retain the meanings explained above, and the value to use for \(\chi\) is the azimuth at the point \(H\) of the arc \(N\).

If the above series are used to integrate the equations
